\documentclass[12pt]{article}

% a template that a friend gave, it's worked well enough for me
% i have added some packages and stuff that have proved useful

\usepackage{fancyhdr}
\usepackage{tipa}
\usepackage{fontspec}
\usepackage{amsfonts}
\usepackage{enumitem}
\usepackage[margin=1in]{geometry}
\usepackage{graphicx}
\usepackage{float}
\usepackage{amsmath}
\usepackage{braket}
\usepackage{amssymb}
\usepackage{booktabs}
\usepackage{hyperref}
\usepackage{mathtools}
\usepackage{xcolor}
\usepackage{float}
\usepackage{algpseudocodex}
\usepackage{titlesec}
\usepackage{bbm}

\pagestyle{fancy}
\fancyhf{} % sets both header and footer to nothing
\lhead{Kevin Sheng}
\setmainfont{Comic Neue}
\renewcommand{\headrulewidth}{1pt}
\setlength{\headheight}{0.75in}
\setlength{\oddsidemargin}{0in}
\setlength{\evensidemargin}{0in}
\setlength{\voffset}{-.5in}
\setlength{\headsep}{10pt}
\setlength{\textwidth}{6.5in}
\setlength{\headwidth}{6.5in}
\setlength{\textheight}{8in}
\renewcommand{\headrulewidth}{0.5pt}
\renewcommand{\footrulewidth}{0.3pt}
\setlength{\textwidth}{6.5in}
\usepackage{setspace}
\usepackage{multicol}
\usepackage{float}
\setlength{\columnsep}{1cm}
\setlength\parindent{24pt}
\usepackage [english]{babel}
\usepackage [autostyle, english = american]{csquotes}
\MakeOuterQuote{"}

\setlength{\parskip}{6pt}
\setlength{\parindent}{0pt}

\titlespacing\section{0pt}{12pt plus 4pt minus 2pt}{0pt plus 2pt minus 2pt}
\titlespacing\subsection{0pt}{12pt plus 4pt minus 2pt}{0pt plus 2pt minus 2pt}
\titlespacing\subsubsection{0pt}{12pt plus 4pt minus 2pt}{0pt plus 2pt minus 2pt}

\hypersetup{colorlinks=true, urlcolor=blue}

\newcommand{\correction}[1]{\textcolor{red}{#1}}


\rhead{Math 168}

\DeclareMathOperator{\Fr}{Fr}
\DeclareMathOperator{\Tr}{Tr}
\newcommand{\lra}{\xLeftrightarrow}
\newcommand{\ra}{\xRightarrow}

\setlist[enumerate,1]{label=(\alph*)}

\begin{document}

\section{Bipartite Networks}

\subsection{Part A}

If a bipartite network has an incidence matrix $B \in \Z^{n \times m}$,
we have nodes $a_1, \cdots, a_n$ and $b_1, \cdots, b_m$, and
there's an edge between $a_i$ and $b_j$ iff $B_{ij}=1$.

I propose that if we number the nodes with $a_1, \cdots, a_n$ coming first
then $b_1, \cdots, b_m$, the adjacency matrix will take the desired form.
Let's formalize this with a map of sorts:
\[v_i=\begin{cases}
            a_i     & i \le n          \\
            b_{i-n} & \text{otherwise}
      \end{cases}\]

Recall that $A_{ij}=1$ iff there's an edge between $v_i$ and $v_j$.
Now we do casework:
\begin{enumerate}[label=\arabic*.]
      \item $i, j \le n$: By the definition of a bipartite graph,
            no two nodes in $a_i, \cdots, a_n$ have edges between them.
            Thus, in this block $A_{ij}=0$ always.
      \item $i \le n, j > n$: This marks the edge $a_i$ to $b_{j-n}$, which is $1$ iff $B_{i(j-n)}=1$.
      \item $i > n, j \le n$: This one is the edge from $b_{i-n}$ to $a_j$, which is set iff $B_{j(i-n)}=1$.
      \item $i, j > n$: Using similar reasoning as in the first case, we get that this block is all $0$s.
\end{enumerate}

\subsection{Part B}

\subsubsection{Forwards Direction}

If a graph is bipartite, we can partition the nodes into sets $A$ and $B$
s.t. all edges are only between these two sets.

I propose that any walk of length $n$ starting in $A$
must end in $A$ if it's of even length and $B$ otherwise.
The proof is by induction, where the base case on a walk of length $0$ is trivial.

For the inductive step on a walk of length $n$ to a walk of length $n+1$, there's two cases:
\begin{itemize}[nolistsep]
      \item If $n$'s even, all walks must end in $A$ and $n+1$ is odd.
            Any additional edge that we take must start in $A$ and end in $B$ by bipartiteness.
            Thus, the hypothesis holds.
      \item If $n$'s odd, all walks must end in $B$ and $n+1$ is even.
            By symmetrical reasoning here, they hypothesis also holds.
\end{itemize}

A cycle is also a walk that starts and ends in the same set.
WLOG let the cycle start in $A$.
If we are to complete the cycle and return to $A$, the cycle must be even-length.
Thus, any bipartite graph must have only even-length cycles.

\subsubsection{Backwards Direction}

I'll assume the graph is connected, since disconnected graph components
can be put into bipartite sets independently.

If any cycle in a graph is even-length, I propose the following partition scheme:
\begin{enumerate}[nolistsep,label=\arabic*.]
      \item Select an arbitrary node $v$ to assign to $A$ first.
      \item All nodes whose shortest path to $v$ is even-length will also be assigned to $A$.
      \item Any remaining nodes (whose shortest path is odd-length) are assigned to $B$.
\end{enumerate}

Now it STP there are no edges between any two nodes in $A$ or $B$.

Assume FSOC there existed an edge between $v_1$ and $v_2$, where $v_1, v_2 \in A$.
Let their shortest path lengths to $v$ be $d_1$ and $d_2$ respectively.
There are two cases.
\begin{enumerate}[nolistsep,label=\arabic*.]
      \item $d_1 < d_2$ or vice versa (WLOG assume the first).
            Since both are even, this means $v_1$'s shortest path is at least $2$
            edges shorter than $v_2$'s.
            This then means that we could've found a smaller $d_2$
            by going to $v_1$ and then taking the $d_1$-length path back to $v$.
      \item $d_1=d_2$, in which case we can construct an odd-length cycle
            by going from $v$ to $v_1$, $v_2$ and back to $v$.
\end{enumerate}
In either case, a contradiction is present.

By analogous reasoning, there are no edges between any two nodes in $B$ as well.

\subsection{Part C}

\section{Linalg With Adjacency Matrices}

\begin{enumerate}
      \item \begin{enumerate}[label=\roman*.]
                  \item $\boxed{A\mathbf{1}}$ gives us $d$, since here $d_i=\sum_{i=1}^{n} A_{ij}$,
                        which indeed adds $1$ to the total for every node connected to node $i$.
                  \item $\boxed{\frac{A\mathbf{1} \cdot \mathbf{1}}{2}}$, since the numerator
                        is the sum of all the degrees and $m$ is half the sum of all the nodes' degrees.
                  \item $\boxed{A^2}$, since for any two nodes one can form a bijection between
                        length-$2$ walks between them and their common neighbors.
            \end{enumerate}
      \item \begin{enumerate}[label=\roman*.]
                  \item $\boxed{A\mathbf{1}}$ in this case gives us the in-degrees,
                        since the $i$th row marks all the other nodes that have edges going \textit{into} $i$.
                  \item $\boxed{A^T\mathbf{1}}$ suffices to give the out-degrees,
                        since transposing the matrix now makes the $i$th row mark all nodes
                        that have edges going into them \textit{from} $i$.
            \end{enumerate}
      \item None of my answers in part a hold for a weighted graph, since all of them rely
            on $A_{ij}=1$ iff there's an edge between $i$ and $j$.
\end{enumerate}

\section{Connectivity}

\subsection{Part A - First Statement}

\subsubsection{Forwards}

If $G$'s connected, any node can be reached from any other node by traversing
some sequence of edges.
Since $G$ itself has $n$ nodes, the maximum length of a path is $n$.

Thus, for any two nodes $i$ and $j$, by the theorem in class $\exists 1 \le k \le n: A^k_{ij} > 0$.
$A$ consists only of nonnegative entries, so this means $\sum_{k=1}^{n} A^k_{ij} > 0$ as well.

\subsubsection{Backwards}

Fix any distinct $i$ and $j$.

$\sum_{k=1}^{n} A^k_{ij} > 0$, and $A$ doesn't have any negative entries.
This means that $\exists 1 \le k \le n: A^k_{ij} > 0$
and there exists a walk of length $k$ between $i$ and $j$.
This is true for all node pairs, so $G$ must be connected. $\square$

\subsection{Part B}

While the statement in part A holds for strong connectivity,
the forward direction breaks down for weak connectivity.
For weakly connected directed graphs, we're only guaranteed that $S+S^T$
has only positive entries.

\section{Newman 6.5}

I solved these in HW 1 already, but I'll put the proofs again here.

\subsection{Part A}

BWOC let there be a 3-regular graph with $2k+1$ nodes.

The number of edges must be both an integer and equal to $\frac{3(2k+1)}{2}$
since every node is connected to three others.

However, these two conditions contradict each other, as $3(2k+1)$ is always odd. $\square$

\subsection{Part B}

A tree with $n$ nodes must have $n-1$ edges.

Since there's $n-1$ edges, the sum of the node degrees is $2(n-1)$.
This makes the average degree $\frac{2(n-1)}{2} < 2$. $\square$

\section{Newman 6.7}

six sevennnnnnnn

\begin{enumerate}
      \item Incoming edges: $\sum_{i=1}^{r} k_i^{in}$

            Outgoing edges: $\sum_{i=1}^{r} k_i^{out}$

      \item If we just use $\sum_{i=1}^{r} k_i^{in}$, we're also counting
            the number of incoming edges that come from nodes $1, 2, \cdots, r$ themselves.

            Thus, we subtract the second expression we got to get
            $\boxed{\sum_{i=1}^{r} k_i^{in}-k_i^{out}}$.

      \item The second inequality is a bit simpler given what we have,
            so I'll prove that first.

            A counting argument should work.
            $\sum_{i=1}^{r} k_i^{in}-k_i^{out}$ accounts for all edges going into nodes $1$ through $r$.

            Any edge coming out from $r+1$ must go into one of $1, 2, \cdots, r$,
            so all edges counted by $k_{r+1}^{out}$ must be also in $\sum_{i=1}^{r} k_i^{in}-k_i^{out}$.
            Thus, $k_{r+1}^{out} \le \sum_{i=1}^{r} k_i^{in}-k_i^{out}$.

            For the first inequality, notice that $\sum_{i=r+1}^{n} k_i^{out}-k_i^{in}$ counts
            the number of edges going out of nodes $r+1, r+2, \cdots, n$.
            It first sums up all edges going out of those nodes unconditionally
            and then subtracts the edges that just go back into $r+1, r+2, \cdots, n$.

            We can use a similar counting argument for this one.
            Any edge that came into node $r$ must have come from nodes $r+1$ through $n$,
            and the RHS counts \textit{all} outgoing edges from that set.
            Thus, $k_r^{in} \le \sum_{i=r+1}^{n} k_i^{out}-k_i^{in}$. $\square$
\end{enumerate}

\section{Triangle Counting}

\begin{enumerate}
      \item There's a 1-6 mapping between triangles in a graph and length-$3$ walks
            that wind up back where the walk started.

            The mapping is as follows.
            If nodes $u$, $v$, and $w$ all share edges, they correspond with these length-$3$ walks:
            \[(u, v, w, u), (u, w, v, u), (v, u, w, v), (v, w, u, v), (w, u, v, w), (w, v, u, w)\]

            Since $A^3_{ij}$ contains the number of length-$3$ walks from $i$ to $j$,
            its trace sums up all length-$3$ walks that start and end at the same position,
            which was exactly the set that had a 1-6 correspondence with the number of triangles.
            Thus, dividing the trace by $6$ gives us the number of triangles.

      \item The formula doesn't work for triangles in a directed graph because
            the mapping between triangles and cycles is reduced to 1-3.
            You can walk through a cycle one way but not necessarily in the other in directed graphs.

      \item $\frac{1}{3}\Tr\left(A^3\right)$ is all we need to make things work.
            The 1-3 mapping turns the triangle $u \to v \to w$ into the following walks:
            \[(u, v, w, u), (v, u, w, v), (w, u, v, w)\]

      \item Each triangle has three nodes, so each triangle adds $1$ to three nodes' $NT_i$.
            To account for this, after we sum up all the $NT_i$s, we need to divide by $3$.

      \item I ran this code:
            \begin{python}
import networkx as nx
import numpy as np

adj = [
    [1],
    [0, 2, 3],
    [1, 4, 5],
    [1, 6],
    [2, 5],
    [2, 4, 6, 7],
    [3, 5, 7],
    [5, 6, 8],
    [7]
]

g = nx.Graph()
a = np.zeros((len(adj), len(adj)), dtype=int)
for i, next_ in enumerate(adj):
    for j in next_:
        if i < j:
            g.add_edge(i, j)
        a[i, j] = 1

tri1 = np.trace(a @ a @ a) // 6
print(f"# of triangles (trace formula): {tri1}")
tri2 = sum(nx.triangles(g).values()) // 3
print(f"# of triangles (nx sum): {tri2}")
            \end{python}
            and got $2$ triangles from both functions, which matches the two triangles
            I found by looking at the graph myself.
\end{enumerate}

\end{document}
